\documentclass[10pt,a4paper]{amsart}
\usepackage[margin=19mm]{geometry}
\usepackage{amsmath,amssymb,mathtools}
\usepackage[scr=rsfs,cal=euler]{mathalfa}
\usepackage{xfrac}
\usepackage[unicode,hidelinks,bookmarks=false]{hyperref}
\newcommand{\ie}{i.e.\ }
\newcommand{\cf}{cf.\ }
\newcommand{\resp}{respectively\ }
\newcommand{\Subsection}{\subsection}
\providecommand{\nobreakdash}{\nobreak\hbox{-}}
\setlength{\emergencystretch}{2em}
\multlinegap0pt
\usepackage[shortlabels]{enumitem}
\usepackage{xcolor}
\usepackage{amssymb}
\newtheorem{lemma}{Lemma}
\newtheorem{theorem}{Theorem}
\usepackage{cleveref}
\crefname{lemma}{Lemma}{Lemmas}
\crefname{theorem}{Theorem}{Theorems}
\newcommand{\normaltt}[1]{\normalfont\texttt{#1}}
\title{Critical exponent of ternary words{ } \\ with few distinct palindromes}
\author{\v{L}ubom\'ira Dvo\v{r}\'akov\'a, Lucas Mol,, Pascal Ochem}
\date{Source: arXiv:2507.04925v3 (CC BY 4.0)}
\begin{document}
\maketitle

\noindent\emph{Source and licence.} Critical exponent of ternary words{ } \\ with few distinct palindromes. Version arXiv:2507.04925v3 (CC BY 4.0), distributed by arXiv under the Creative Commons Attribution 4.0 International licence, http://creativecommons.org/licenses/by/4.0/, as recorded in the arXiv licence metadata for this exact version and shown on the abstract page https://arxiv.org/abs/2507.04925v3. Full licence notice: \texttt{licenses/arxiv-2507.04925-CC-BY-4.0.txt}.\par\smallskip

\noindent\textit{Editorial scope.} Counted unit: the article's theorem 41\_22 (source0.tex lines 700--706). The article's complete original proof of it is delivered with it (source0.tex lines 707--755, one \texttt{proof} environment of 49 lines). No mathematics is written, repaired or re-derived: the statement and the proof are the article's own lines, in the article's own notation, and no retained line is substituted or re-flowed.  Retained uncounted as the source-local context the proof needs: lines 637--639; lines 821--838.  Nothing was omitted from any retained span, so the package carries no omission record.  No retained line contains a citation, so the wrapper carries no bibliography and no bibliography entry.  No retained line contains a figure environment, an image-inclusion command or drawing code, so no image is delivered and the package declares no asset dependency.  Editorial formatting only, in addition: an AMS \texttt{amsart} preamble that restates the article's own declarations and macros used by the retained lines, namely the environments \texttt{lemma}, \texttt{theorem} on separate counters, together with the standard packages those lines need.  The retained environments print in this document as Lemma 1, Theorem 1 in order of appearance, whereas the article prints the same items with the numbering of its own full document.  The complete native source of the article as distributed in the arXiv e-print for this exact version is bundled unchanged as \texttt{sources/181139/source0.tex}.
\begin{lemma}\label{ggood}
A bi-infinite word is $\gamma$-good if and only if it has the same factor set as $\gamma(\eta^\omega(\normaltt{0}))$.
\end{lemma}

% \begin{align*}
%     g'''(\texttt{0}) &= \texttt{01202120121}\\
%     g'''(\texttt{1}) &= \texttt{0120102}\\
%     g'''(\texttt{2}) &= \texttt{012101202120102}\\
%     g'''(\texttt{3}) &= \texttt{012021201020121}
% \end{align*}
% \end{minipage}

% Unfortunately, the fixed point of $h'$ ( resp. $h''$, $h'''$) does not seem to have a characterization.

\section{Critical exponent of \texorpdfstring{$g(h^{\omega}({\texttt 0}))$}{g(h omega(0))}}\label{sec:ce}
Let us recall the definition of the morphisms $h$ and $g$:
\begin{align*}
    h(\texttt{0}) &= \texttt{01213012} & g(\texttt{0}) &= \texttt{0102012}\\
    h(\texttt{1}) &= \texttt{31} & g(\texttt{1}) &= \texttt{0212}\\
    h(\texttt{2}) &= \texttt{01201312} & g(\texttt{2}) &= \texttt{0121012}\\
    h(\texttt{3}) &= \texttt{0121301312} & g(\texttt{3}) &= \texttt{01020121012}
\end{align*}

\begin{theorem}\label{41_22}{\ }
\begin{enumerate}[label=\normalfont(\alph*)]
\item $\gamma(\eta^\omega(\normaltt{0}))$ contains exactly 16 palindromes and is $\tfrac{41}{22}^+$-free.\label{it:16p}
\item Every bi-infinite ternary $\tfrac{52}{27}$-free word containing at most $16$ palindromes has the same factor set as either $\gamma(\eta^\omega(\normaltt{0}))$ or $\gamma(\eta^\omega(\normaltt{0}))^R$.\label{it:pos}
\item Every recurrent ternary $\tfrac{25}{13}$-free word containing at most $17$ palindromes has the same factor set as either $\gamma(\eta^\omega(\normaltt{0}))$ or $\gamma(\eta^\omega(\normaltt{0}))^R$.\label{it:neg}
\end{enumerate}
\end{theorem}

\begin{proof}
Let us prove \Cref{it:16p}.  The word
$\gamma(\eta^\omega(\texttt{0}))$ contains the palindromes $\varepsilon$, \texttt{0}, \texttt{1}, \texttt{2}, the six palindromes
of the form $bcb$, and the six palindromes of the form $abcba$. It is easy to check that it contains no other palindrome, i.e., no $cabcbac$.
To show that $\gamma(\eta^\omega(\texttt{0}))$ is $\tfrac{41}{22}^+$-free, we consider the morphic word
$g(h^\omega(\texttt{0}))$ defined by the following morphisms $h$ and $g$.
\begin{align*}
    h(\texttt{0}) &= \texttt{01213012} & g(\texttt{0}) &= \texttt{0102012}\\
    h(\texttt{1}) &= \texttt{31} & g(\texttt{1}) &= \texttt{0212}\\
    h(\texttt{2}) &= \texttt{01201312} & g(\texttt{2}) &= \texttt{0121012}\\
    h(\texttt{3}) &= \texttt{0121301312} & g(\texttt{3}) &= \texttt{01020121012}
\end{align*}
We check that $g(h^\omega(\texttt{0}))$ avoids $F_\gamma$. We will show in Section~\ref{sec:ce} that $g(h^\omega(\texttt{0}))$ is $\tfrac{41}{22}^+$-free.
Since $g(h^\omega(\texttt{0}))$ is square-free and avoids $F_\gamma$, we see that $g(h^\omega(\texttt{0}))$ is $\gamma$-good. So by Lemma~\ref{ggood}, the words
$g(h^\omega(\texttt{0}))$ and $\gamma(\eta^\omega(\texttt{0}))$ have the same factor set.\footnote{Notice that $g(h^\omega(\texttt{0}))$ is not bi-infinite,
as required by Lemma~\ref{ggood}. However, $g(h^\omega(\texttt{0}))$ is uniformly recurrent, which is somehow a stronger condition for our purpose.}
Thus, $\gamma(\eta^\omega(\texttt{0}))$ is also $\tfrac{41}{22}^+$-free.

Let us prove \Cref{it:pos}.
We say that a ternary word (finite or infinite) is 16-good if it is $\tfrac{52}{27}$-free and contains at most $16$ palindromes.
We construct the set $S_{16}^{36}$ defined as follows:
a word $v$ is in $S_{16}^{36}$ if and only if there exists a 16-good word $pvs$ such that $|p|=|v|=|s|=36$.
Then we construct the Rauzy graph $G_{16}^{36}$ based on $S_{16}^{36}$, that is, such that vertices correspond to factors of length 35 and arcs correspond to factors of length 36.
We notice that $G_{16}^{36}$ consists of two connected components that are images of each other with respect to reversal.
Moreover, the connected component containing the factor \texttt{0120} is equal to the Rauzy graph
of the factors of length 36 of $\gamma(\eta^\omega(\texttt{0}))$.
Let $w$ be a bi-infinite 16-good word belonging to this connected component, that is, $w$ contains \texttt{0120}.
Since $\max_{f\in F_\gamma}|f|=36$, we see that $w$ is $\gamma$-good. By Lemma~\ref{ggood}, we conclude that $w$ has the same factor set as $\gamma(\eta^\omega(\texttt{0}))$.
By symmetry, this proves \Cref{it:pos}.

Let us prove \Cref{it:neg}.
We say that a ternary word (finite or infinite) is 17-good if it is $\tfrac{25}{13}$-free and contains at most $17$ palindromes.
Let $w$ be a recurrent 17-good word.
We check by backtracking that no infinite 17-good word avoids $\texttt{010}$.
By symmetry, $w$ must contain the six palindromes of the form $bcb$ and thus $abcba$.
So $w$ contains the 16 palindromes of $\gamma(\eta^\omega(\texttt{0}))$.
Notice that $w$ cannot contain a palindrome of the form $abacaba$ since $abacaba$
cannot be extended into a ternary square-free word.
So, without loss of generality, we assume that the potential 17th palindrome of $w$ is \texttt{0120210}.
We say that a ternary word (finite or infinite) is 17-great if it is 17-good and every palindrome it contains is a factor of \texttt{0120210}, \texttt{01210}, \texttt{02120}, \texttt{10201}, \texttt{20102}, or \texttt{21012}.
Thus $w$ is 17-great. We construct the set $S_{17}^{36}$ defined as follows:
a word $v$ is in $S_{17}^{36}$ if and only if there exists a 17-great word $pvs$ such that $|p|=|v|=|s|=36$.
Then we construct the Rauzy graph $G_{17}^{36}$ based on $S_{17}^{36}$.
Since $w$ is recurrent, $w$ is in a strongly connected component of $G_{17}^{36}$.
We have checked that the graph induced by each strongly connected component of $G_{17}^{36}$
is isomorphic to $G_{16}^{36}$.
Similarly, this shows that $w$ has the same factor set as either
$\gamma(\eta^\omega(\texttt{0}))$ or $\gamma(\eta^\omega(\texttt{0}))^R$.
\end{proof}

\end{document}
